% ============================================================================
% Fichier  : partie3/P03_C06_loi_2024.tex
% Titre    : La Loi n° 2024/017 : Autorisation, Autorité, Sanctions
% Auteur   : MINKA MI NGUIDJOÏ Thierry Emmanuel
% Version  : 3.0 -- novembre 2026
% Corrigés : corriges/Corr_C06.tex
% ============================================================================
% Sources : loi n° 2024/017 (texte officiel, titres I à VII) ; loi
% n° 2010/012, art. 41 à 47 (texte officiel) ; Convention de Malabo ;
% Convention de Budapest ; acte additionnel CEDEAO A/SA.1/01/10 ; RGPD.
% Statut des conventions d'après Minka Mi Nguidjoï, Massoussi Songo et
% Bialo (2026, manuscrit) : Budapest, Cameroun Partie le 15/12/2023, en
% vigueur le 01/04/2024 ; Malabo, signée le 12/08/2021, non ratifiée.
% ============================================================================

\chapter{La Loi n° 2024/017 : Autorisation, Autorité, Sanctions}
\label{chap:loi}
\chaptermark{La loi n° 2024/017}

\epigraph{%
    « On cite une loi pour ce qu'elle dit,\\
    pas pour ce que l'on aimerait qu'elle dise. »%
}{Principe du cours}

\niveauChapitre{Fondamental}{%
    Chapitres~\ref{chap:donnees}, \ref{chap:principes} et
    \ref{chap:international} à~\ref{chap:normes}~; droit du numérique
    (UE MSN-2S3-JUR01)
}

\begin{boxobjectifs}
\begin{itemize}
    \item Présenter l'architecture de la loi n°~2024/017 en sept titres.
    \item Expliquer le régime d'autorisation préalable et ses conséquences
          tant que l'Autorité n'est pas organisée.
    \item Décrire les missions de l'Autorité et l'échelle des sanctions
          administratives, civiles et pénales.
    \item Dire ce que la loi a internalisé des sources internationales et
          ce qu'elle a ajouté en propre.
    \item Recenser les dispositions qui attendent un texte d'application et
          en tirer les conséquences pour la cotation.
    \item Lire les dispositions transitoires et la clause d'abrogation.
\end{itemize}
\end{boxobjectifs}

\bigskip

% ============================================================================
\section*{Avant de commencer : « où dépose-t-on la demande ? »}
% ============================================================================

Bafoussam, jeudi 25~juin~2026. Le lendemain de la lettre de M.~TAGNE, Me~ETOUNDI
lit la loi au directeur général, article par article. À l'article~19,
M.~TCHOUA l'arrête~: « Une autorisation pour chaque traitement. Très bien.
Où dépose-t-on la demande~? »

Me~ETOUNDI tourne les pages jusqu'à l'article~53. L'Autorité de protection
existe dans la loi. Sa création, son organisation et son fonctionnement
relèvent d'un décret qu'elle n'a pas trouvé. Les modalités de l'autorisation
relèvent d'un autre texte. « Je ne sais pas, répond-elle. Mais je sais ce
qu'il y aura dans le dossier le jour où le guichet ouvrira. »

\begin{boxalert}
\textbf{Violation typique.} Conclure que la loi ne s'applique pas encore
parce que l'Autorité n'est pas organisée, et attendre.

\smallskip
\textbf{Règle~:} la loi est en vigueur et son délai de mise en conformité est
expiré. Les obligations qui ne dépendent d'aucun texte d'application
s'imposent dès maintenant. Pour les autres, l'organisation prépare ce qui
dépend d'elle et l'auditeur cote séparément l'exigence et la diligence.
\end{boxalert}

% ============================================================================
\section{Ce que la Loi a Internalisé}
\label{sec:c06-internalise}
% ============================================================================

Les chapitres précédents ont décrit le paysage international. La loi
n°~2024/017 en est l'héritière sans en être la copie. Le tableau
ci-dessous sépare ce qu'elle reprend de ce qu'elle ajoute.

\begin{tableaucours}{L{2.8cm}L{4.4cm}L{4.8cm}}{%
    \tableauHeader{Titre de la loi} &
    \tableauHeader{Ce qu'elle reprend} &
    \tableauHeader{Ce qu'elle ajoute ou change}
}
I. Définitions, champ & Vocabulaire commun à la Convention~108, à Malabo et
au RGPD & Protection des personnes « en transit »~; liste ouverte de données
sensibles, dont les transactions bancaires et l'origine régionale \\
II. Principes & Les huit principes de l'OCDE & Consentement comme règle,
trois dérogations seulement~; pas de principe général de minimisation \\
II. Formalités & Formalités préalables de Malabo & Autorisation pour tout
traitement, sans régime de déclaration \\
II. Obligations & Registre, évaluation d'impact, sous-traitance, transferts
du RGPD & Mêmes obligations pour le sous-traitant~; convention publique~;
quatre vérifications cumulées pour un transfert~; rapport annuel de
sécurité \\
III. Droits & Droits du RGPD, portabilité comprise & Directives après le
décès~; délais renvoyés au règlement \\
V. Autorité & Autorité indépendante des principes de l'ONU et de Malabo &
Création et organisation par décret \\
VI. Sanctions & Amendes administratives & Délits nombreux, dont le
profilage~; responsabilité pénale des personnes morales \\
\end{tableaucours}
\vspace{-0.8em}
\begin{center}{\small\itshape Tableau~6.1. Ce que la loi reprend et ce
qu'elle ajoute}\end{center}

La colonne de droite est celle que l'auditeur doit connaître par cœur. Elle
contient tout ce qu'un consultant formé au seul RGPD oubliera
(Chapitre~\ref{chap:rgpd}).

% ============================================================================
\section{L'Architecture de la Loi}
\label{sec:c06-architecture}
% ============================================================================

\begin{tableaucours}{cL{5.6cm}cL{4.0cm}}{%
    \tableauHeader{Titre} &
    \tableauHeader{Objet} &
    \tableauHeader{Art.} &
    \tableauHeader{Chapitres du cours}
}
I & Dispositions générales~: objet, champ, définitions & 1--5 &
\ref{chap:donnees} \\
II & Régime juridique~: principes, formalités préalables, obligations,
interconnexion & 6--36 & \ref{chap:principes}, \ref{chap:cartographie},
\ref{chap:risque}, \ref{chap:securite}, \ref{chap:soustraitants},
\ref{chap:transferts}, \ref{chap:violations} \\
III & Droits de la personne concernée & 37--47 & \ref{chap:droits},
\ref{chap:profilage} \\
IV & Interdictions & 48--52 & \ref{chap:donnees}, \ref{chap:secteurs} \\
V & Autorité de protection & 53 & ce chapitre \\
VI & Sanctions administratives, civiles et pénales & 54--71 & ce chapitre \\
VII & Dispositions diverses, transitoires et finales & 72--75 &
\ref{chap:prologue}, \ref{chap:secteurs} \\
\end{tableaucours}
\vspace{-0.8em}
\begin{center}{\small\itshape Tableau~6.2. Les sept titres de la loi n°~2024/017}\end{center}

La loi compte soixante-quinze articles. Elle « abroge toutes les
dispositions antérieures contraires » (article~75) et renvoie à des textes
particuliers les modalités d'application (article~74). Une quinzaine de ses
dispositions dépendent d'un texte réglementaire ou d'un référentiel de
l'Autorité~: critères de risque, évaluation d'impact, durées de
conservation, délais des droits, modalités d'autorisation et de contrôle.

% ============================================================================
\section{Ce qui Attend un Texte d'Application}
\label{sec:c06-application}
% ============================================================================

\begin{tableaucours}{L{4.6cm}cL{5.6cm}}{%
    \tableauHeader{Objet} &
    \tableauHeader{Art.} &
    \tableauHeader{Conséquence pour l'audit}
}
Création et organisation de l'Autorité & 53 & Pas de guichet, pas de
contrôle administratif \\
Modalités de l'autorisation & 19, 32 & Double cote~: exigence et diligence \\
Critères et modalités de l'évaluation du risque & 25 & Méthode déclarée par
l'organisation \\
Évaluation d'impact & 33 & Critères de risque élevé déclarés \\
Durées de conservation & 13, 28 & Durées fixées et justifiées par
l'organisation \\
Délais d'examen des demandes & 47 & Critère de substitution déclaré \\
Référentiel des mesures & 24, 27, 57 & Contenu du rapport annuel inconnu \\
Mécanisme de certification & 34, 59 & Sanction de l'article~59 sans objet
pour l'instant \\
Modalités du contrôle & 35 & Sans objet pour l'instant \\
Liste des pays, clauses types & 32, 53 & Deux vérifications de l'article~32
non cotables \\
Textes des secteurs & 4, 72 & Loi générale appliquée, réserve signalée \\
\end{tableaucours}
\vspace{-0.8em}
\begin{center}{\small\itshape Tableau~6.3. Les dispositions suspendues à un
texte (état en novembre~2026, tableau à refaire à chaque mission)}\end{center}

Chaque ligne de ce tableau produira dans les chapitres suivants un \NEK{} de
critère, c'est-à-dire une incertitude qui tient au droit et non à
l'organisation (Chapitre~\ref{chap:audit}). L'auditeur refait le tableau au
début de chaque mission. Un texte publié entre-temps transforme un \NEK{} en
exigence cotable.

% ============================================================================
\section{Le Régime d'Autorisation Préalable}
\label{sec:c06-autorisation}
% ============================================================================

\begin{boxjuris}
\textbf{\loicam{n°~2024/017}, article~19}~\cite{loi2024017}~:
\textit{« (1) Le traitement des données à caractère personnel est soumis à
l'obtention préalable d'une autorisation délivrée par l'Autorité de
protection des données à caractère personnel. (2) Tout processus
d'interconnexion et d'interopérabilité des fichiers de données sensibles,
relatives aux mineurs, est soumis à l'obtention préalable d'une
autorisation délivrée par l'Autorité de protection des données à caractère
personnel. (3) Les modalités de délivrance de l'autorisation visée aux
alinéas 1 et 2 ci-dessus sont fixées par voie réglementaire. »}

\smallskip
\textbf{Portée.} La loi ne prévoit ni déclaration simplifiée ni dispense~:
tout traitement est soumis à autorisation. L'article~49 interdit le
traitement sans autorisation préalable, et l'article~55 le punit d'une
pénalité de 5 à 50~millions de francs CFA.
\end{boxjuris}

Le régime est l'inverse du modèle européen, qui a supprimé les formalités
préalables au profit de la responsabilisation du responsable. Il rapproche
la loi camerounaise des premières lois africaines et françaises, fondées
sur un contrôle \textit{a priori}.

Ce régime pose une difficulté pratique~: l'autorisation est délivrée par une
Autorité dont « la création, ainsi que l'organisation et le
fonctionnement » sont fixés par décret (article~53, alinéa~2), et selon des
modalités elles-mêmes réglementaires (article~19, alinéa~3). Tant que ces
textes ne sont pas publiés et que l'Autorité ne reçoit pas de demandes,
l'organisation ne peut pas obtenir d'autorisation. Que doit coter
l'auditeur~?

\begin{itemize}
    \item « Le traitement est autorisé »~: \NON{}, puisqu'aucune autorisation
          n'existe. La cote est exacte~; elle ne dit pas que l'écart est
          imputable à l'organisation.
    \item « L'organisation a fait ce qui dépendait d'elle »~: \OUI{} si elle
          a préparé le dossier (registre, évaluations, mesures) et peut
          prouver qu'elle le déposera dès que la procédure sera ouverte~;
          \NON{} sinon.
\end{itemize}

Le rapport présente les deux cotes ensemble. À NYANGA, la seconde est
\NON{}~: aucun dossier n'a été préparé.

% ============================================================================
\section{L'Autorité de Protection}
\label{sec:c06-autorite}
% ============================================================================

\begin{boxjuris}
\textbf{\loicam{n°~2024/017}, article~53, alinéa~1 (extraits)}~\cite{loi2024017}~:
\textit{« L'Autorité de protection des données à caractère personnel est un
organisme public indépendant chargé, notamment~: de veiller à l'application
de la présente loi, des textes pris pour son application ainsi que des
conventions internationales en la matière~; de délivrer les autorisations
prévues par la présente loi, assorties des cahiers de charges
conséquents~; d'élaborer, publier et mettre à jour le référentiel des
mesures techniques et organisationnelles des données à caractère
personnel~; [\ldots] de traiter les réclamations, pétitions et plaintes
[\ldots]~; d'établir et publier la liste des pays reconnus comme offrant un
niveau de protection des données à caractère personnel équivalent aux
exigences prévues par le droit camerounais~; [\ldots] »}
\end{boxjuris}

Trois instruments de l'Autorité conditionnent l'audit~: le
\textbf{référentiel des mesures} (articles~24, 27, 28, 57), les
\textbf{cahiers des charges} attachés aux autorisations (article~61), et la
\textbf{liste des pays} offrant une protection équivalente
(Chapitre~\ref{chap:transferts}). Tant qu'ils n'existent pas, les exigences
qui en dépendent sont des \NEK{} de critère.

\begin{boxinfo}
\textbf{État de l'Autorité au moment de la rédaction.} À la date de ce cours
(novembre~2026), l'auteur n'a pas trouvé de texte publié organisant
l'Autorité, et ses coauteurs et lui relèvent dans leurs travaux en cours que
ce texte était attendu après l'échéance du 23~juin~2026. Le lecteur
vérifiera la situation au moment où il lit~: si l'Autorité est organisée,
une partie des \NEK{} de critère de ce livre se lève.
\end{boxinfo}

% ============================================================================
\section{Contrôle et Certification}
\label{sec:c06-certification}
% ============================================================================

L'article~35 confie à l'Autorité le contrôle des traitements, selon des
modalités réglementaires. L'article~34 lui demande en outre de mettre en
place une certification.

\begin{boxjuris}
\textbf{\loicam{n°~2024/017}, article~34}~\cite{loi2024017}~:
\textit{« L'Autorité de protection des données à caractère personnel met en
place un mécanisme de certification du traitement des données à caractère
personnel conforme aux principes et exigences énoncés par la présente
loi. »}

\smallskip
\textbf{Portée.} L'article~5 définit la certification comme un outil
permettant de communiquer sur le niveau de protection « au regard des
critères d'un référentiel approuvé au préalable par l'Autorité ». Il s'agit
d'une certification de conformité à la loi, distincte des certifications
ISO. L'article~59 punit le responsable qui « opère sans certification
prévue dans la présente loi »~: tant que le mécanisme n'existe pas, la
portée de cette sanction est un \NEK{} de critère.
\end{boxjuris}

\begin{tableaucours}{L{3.2cm}L{4.4cm}L{4.4cm}}{%
    \tableauHeader{} &
    \tableauHeader{Certification de l'article~34} &
    \tableauHeader{Certification ISO/IEC 27701}
}
Délivrée par & L'Autorité, ou selon son mécanisme & Un organisme de
certification accrédité \\
Référentiel & Approuvé par l'Autorité & Norme internationale \\
Objet & Traitement, processus, produit & Système de management sur un
périmètre \\
Portée juridique & Prévue par la loi (articles~34, 59) & Aucune au regard de
la loi \\
État en novembre~2026 & Mécanisme non publié & Disponible \\
\end{tableaucours}
\vspace{-0.8em}
\begin{center}{\small\itshape Tableau~6.4. Deux certifications à ne pas confondre (voir aussi le Chapitre~\ref{chap:normes})}\end{center}

% ============================================================================
\section{Les Sanctions}
\label{sec:c06-sanctions}
% ============================================================================

\subsection{Sanctions administratives}

\begin{boxjuris}
\textbf{\loicam{n°~2024/017}, article~54}~\cite{loi2024017}~:
\textit{« (1) Lorsque l'Autorité de protection des données à caractère
personnel constate le non-respect des obligations mises à la charge du
responsable du traitement ou de son sous-traitant, celui-ci est mis en
demeure de s'y conformer dans un délai maximum de dix (10) jours. (2) Passé
le délai prévu à l'alinéa 1 ci-dessus, l'Autorité de protection des données
à caractère personnel adresse une injonction de mise en conformité du
traitement, sous astreinte dont le montant ne peut excéder cent mille
(100~000) francs CFA par jour de retard. [\ldots] »}

\smallskip
\textbf{Portée.} À défaut de mise en conformité, l'Autorité peut suspendre
l'activité, retirer l'autorisation ou interdire tout traitement
(alinéa~3).
\end{boxjuris}

\begin{tableaucours}{cL{6.4cm}L{4.0cm}}{%
    \tableauHeader{Art.} &
    \tableauHeader{Manquement} &
    \tableauHeader{Montant (F CFA)}
}
55 & Traitement sans autorisation préalable & 5 à 50~millions \\
56 & Refus de communiquer à la personne les informations la concernant & 1
à 10~millions \\
57 & Non-respect du référentiel & 5 à 20~millions \\
58 & Non-respect des règles d'interconnexion & 1 à 5~millions \\
59 & Traitement sans la certification prévue & 5 à 15~millions \\
60 & Transfert vers un pays tiers sans autorisation & 10 à 50~millions \\
61 & Non-respect du cahier des charges & 10 à 100~millions \\
\end{tableaucours}
\vspace{-0.8em}
\begin{center}{\small\itshape Tableau~6.5. Pénalités et amendes administratives (articles~55 à~61)}\end{center}

\subsection{Sanctions civiles et pénales}

L'article~62 ouvre à la personne concernée le référé « en cas d'atteinte
grave aux droits », éventuellement sous astreinte, et l'action en
réparation. Les articles~63 à~70 créent des délits~; l'article~71 rend les
personnes morales pénalement responsables, avec une amende de 50~millions à
un milliard de francs CFA.

\begin{tableaucours}{cL{6.6cm}L{3.8cm}}{%
    \tableauHeader{Art.} &
    \tableauHeader{Délit} &
    \tableauHeader{Peine d'emprisonnement}
}
63 & Collecte ou accès frauduleux, déloyal ou illicite (peines doublées
avec verrouillage ou cryptage) & 2 à 5~ans \\
64 & Traitement maintenu malgré l'opposition, refus de rectifier ou
d'effacer & 1 à 3~ans \\
65 & Traitement à des fins de profilage & 3 à 10~ans \\
66 & Données sensibles sans consentement exprès~; récupération
frauduleuse de données supprimées & 6~mois à 2~ans \\
67 & Détournement de finalité, traitement ultérieur incompatible & 6~mois
à 2~ans \\
68 & Divulgation portant atteinte à la dignité ou à la vie privée (sur
plainte) & 6~mois à 2~ans \\
69 & Transfert vers l'étranger en méconnaissance de la loi & 3 à 10~ans \\
70 & Entrave à l'action de l'Autorité & 6~mois à 3~ans \\
\end{tableaucours}
\vspace{-0.8em}
\begin{center}{\small\itshape Tableau~6.6. Les délits de la loi (amendes assorties, non reproduites)}\end{center}

L'article~65 mérite attention~: il punit « le responsable de traitement ou
le sous-traitant qui procède ou fait procéder à un traitement de données à
caractère personnel à des fins de profilage », sans condition apparente.
Sa portée exacte, pour une notation de crédit par exemple, est discutée au
Chapitre~\ref{chap:profilage}.

% ============================================================================
\section{Dispositions Transitoires et Abrogation}
\label{sec:c06-transitoires}
% ============================================================================

Trois articles ferment la loi. L'article~73 a donné dix-huit mois aux
organisations pour se conformer~; ce délai a expiré le 23~juin~2026
(Chapitre~\ref{chap:prologue}). L'article~74 renvoie les modalités
d'application à des textes particuliers. L'article~75 abroge « toutes les
dispositions antérieures contraires ».

La formule de l'article~75 ne nomme aucun texte. Elle laisse donc en vigueur
tout ce qui, dans les lois antérieures, n'est pas contraire à la loi
nouvelle. La loi n°~2010/012 relative à la cybersécurité et à la
cybercriminalité contient une quinzaine d'articles sur la vie privée, la
conservation des données et les infractions qui s'y rapportent. Lesquels
subsistent~? Aucune liste officielle ne le dit. C'est à l'auditeur et au
juriste de le déterminer article par article, et c'est l'objet du
Chapitre~\ref{chap:loi2010}.

% ============================================================================
\section{Ce que la Loi demande à chaque Métier}
\label{sec:c06-metiers}
% ============================================================================

\begin{boxmetier}{Juriste conformité}
Lire le texte officiel, distinguer les obligations applicables, les
obligations suspendues à un texte d'application et les bonnes pratiques
importées, et le dire dans chaque avis.
\end{boxmetier}

\begin{boxmetier}{Auditeur de la protection des données}
Coter séparément « l'exigence est satisfaite » et « l'organisation a fait ce
qui dépendait d'elle » quand l'écart tient à l'absence de l'Autorité.
\end{boxmetier}

\begin{boxmetier}{Responsable de traitement}
Préparer dès maintenant les dossiers d'autorisation~: registre, fondements,
évaluations, mesures. Le jour où la procédure ouvrira, ce sont eux qui
seront demandés.
\end{boxmetier}

\begin{boxresume}
\begin{itemize}
    \item La loi compte sept titres et soixante-quinze articles~; une
          quinzaine de dispositions attendent un texte d'application.
    \item Tout traitement est soumis à autorisation préalable de
          l'Autorité (article~19)~; tant qu'elle n'est pas organisée,
          l'auditeur cote aussi ce qui dépendait de l'organisation.
    \item Les sanctions vont de la mise en demeure à dix jours jusqu'à un
          milliard de francs CFA pour une personne morale.
    \item La loi internalise les principes internationaux et y ajoute
          l'autorisation préalable, un consentement de principe, une liste
          large de données sensibles et de nombreux délits.
    \item L'article~75 n'abroge que les dispositions contraires~: la loi
          n°~2010/012 subsiste pour le reste.
\end{itemize}
\end{boxresume}

\begin{boxinfo}
\textbf{Points de contrôle.} Vous avez compris ce chapitre si vous savez~:
\begin{itemize}
    \item citer trois ajouts de la loi par rapport à ses sources~;
    \item expliquer comment coter l'absence d'autorisation quand l'Autorité
          n'est pas organisée~;
    \item rattacher un écart à l'article qui le sanctionne.
\end{itemize}
\end{boxinfo}

\begin{boxexo}[title={\ding{45}\quad Exercice~6.1. Quel article \niveaubase}]
Indiquez l'article qui sanctionne~: (a) un transfert vers un hébergeur
étranger sans autorisation~; (b) un agent qui récupère des fichiers
effacés de la corbeille d'un collègue~; (c) un établissement qui refuse de
communiquer à un client ses données~; (d) une entreprise qui refuse de
recevoir les agents de l'Autorité.
\end{boxexo}

\begin{boxexo}[title={\ding{45}\quad Exercice~6.2. Sans Autorité \niveauexpert}]
Un client de MBOLO dépose plainte auprès de l'Autorité en septembre~2026.
L'Autorité n'est pas encore organisée. (a) Quelles autres voies s'offrent
à lui~? (b) Quel est l'intérêt, pour MBOLO, d'avoir préparé ses dossiers
d'autorisation~? (c) Comment l'auditeur rend-il compte de cette
situation dans son rapport~?
\end{boxexo}

\begin{boxexo}[title={\ding{45}\quad Exercice~6.3. Reprise ou ajout \niveaumoyen}]
Pour chaque disposition, dites si la loi reprend une source internationale
(laquelle) ou innove~: (a) la portabilité de l'article~43~; (b) les
transactions bancaires parmi les données sensibles~; (c) l'autorisation
préalable de l'article~19~; (d) la convention publique de l'article~31~;
(e) le droit d'accès de l'article~39.
\end{boxexo}

\begin{boxexo}[title={\ding{45}\quad Exercice~6.4. Le tableau des textes attendus \niveauexpert}]
Un décret organisant l'Autorité paraît pendant votre mission.
\begin{enumerate}[(a)]
    \item Quelles lignes du Tableau~6.3 changent, et lesquelles ne changent
          pas~?
    \item Quelles cotes \NEK{} de vos papiers de travail devez-vous
          reprendre~?
    \item Rédigez le paragraphe du rapport qui explique au commanditaire
          pourquoi des cotes ont changé sans que l'organisation ait rien
          fait.
\end{enumerate}
\end{boxexo}

\bigskip

\begin{boxinfo}
\textbf{Pour aller plus loin.}
\begin{itemize}
    \item La loi n°~2024/017~\cite{loi2024017}, dans son texte officiel,
          titres~V à~VII.
    \item La Convention de Malabo~\cite{ua2014malabo}, chapitre~II, pour
          mesurer ce que la loi en a repris.
\end{itemize}
\end{boxinfo}

\begin{boxinfo}
\textbf{Transition.} La loi de 2024 n'est pas seule. Une loi de 2010, écrite
pour la sécurité des réseaux, impose encore de conserver, de protéger et de
ne pas intercepter. C'est l'objet du Chapitre~\ref{chap:loi2010}.
\end{boxinfo}
